\SourceLeaf{176}%
% A544-P0176-D01
noch keinesweges dazu bestimmt, den überaus reichhaltigen Gegenstand zu erschöpfen. Die Entwickelung der \emph{allgemeinen} Theorie, welche eine ganz eigenthümliche Erweiterung des Feldes der höhern Arithmetik erfordert, bleibt vielmehr der künftigen Fortsetzung vorbehalten, während in diese erste Abhandlung diejenigen Untersuchungen aufgenommen sind, welche sich ohne eine solche Erweiterung vollständig darstellen liessen. Von den Resultaten kann in dieser Anzeige nur ein Theil ausgehoben werden.

% A544-P0176-D02
Eine ganze Zahl $a$ heisst biquadratischer Rest der ganzen Zahl $p$, wenn es Zahlen der Form $x^4-a$ gibt, die durch $p$ theilbar sind; biquadratischer Nichtrest hingegen, wenn keine Zahlen jener Form durch $p$ theilbar sein können. Offenbar sind alle biquadratischen Reste von $p$ zugleich quadratische Reste derselben Zahl, und also alle quadratischen Nichtreste auch biquadratische Nichtreste: allein nicht alle quadratischen Reste sind zugleich biquadratische Reste. Es ist zureichend, die Untersuchungen auf den Fall einzuschränken, wo $p$ eine Primzahl von der Form $4n+1$, und $a$ nicht durch $p$ theilbar ist, da alle anderen Fälle sich leicht auf diesen zurückführen lassen.

% A544-P0176-D03
Die Untersuchungen über diesen Gegenstand zerfallen in zwei Abtheilungen, je nachdem $p$ oder $a$ als gegeben angesehen wird. Die erstere ist von viel geringerer Schwierigkeit als die zweite, und verglichen mit letzterer als ganz elementarisch zu betrachten. Alles Wesentliche, was darüber zu sagen ist, enthält die Abhandlung vollständig.

% A544-P0176-D04
Aus der zweiten Abtheilung hingegen sind hier nur erst einige specielle Fälle abgehandelt, die sich ohne zu grosse Zurüstungen abmachen liessen, und als Vorbereitungen zu der künftig zu gebenden allgemeinen Theorie dienen können. Dies sind diejenigen, wo $a=-1$, und $a=\pm2$ gesetzt wird. Der erstere Fall hat gar keine Schwierigkeit: es war auch schon in dem Werke, \emph{Disquisitiones Arithmeticae}, gezeigt, dass $-1$ ein biquadratischer Rest von $p$ ist, so oft $p$ die Form $8n+1$ hat, hingegen ein blos quadratischer Rest und biquadratischer Nichtrest von $p$, wenn $p$ von der Form $8n+5$ wird. Ganz anders verhält es sich mit dem Fall $a=\pm2$. Es ist zwar längst bekannt, dass $+2$ und $-2$ von $p$ quadratische und also auch biquadratische Nichtreste sind, wenn $p$ die Form $8n+5$ hat, und wenigstens quadratische Reste, wenn $p$ von der Form $8n+1$ ist, wie auch dass bei dieser Form von $p$ entweder $+2$ und $-2$ zugleich biquadratische Reste, oder zugleich biquadratische Nichtreste werden: 
\SourceLeaf{177}%
% A544-P0177-D01
allein die Unterscheidung, welcher dieser beiden Fälle eintrete, ist eine Untersuchung von viel höherer Art, und es werden dazu in der Abhandlung zwei verschiedene Criterien entwickelt.

% A544-P0177-D02
Das erste Criterium hängt mit der Zerlegung der Zahl $p$ in ein einfaches und ein doppeltes Quadrat zusammen, die bekanntlich (da, wie schon bemerkt ist, angenommen wird, dass $p$ eine Primzahl sei) immer möglich und nur auf Eine Art möglich ist. Setzt man $p=gg+2hh$, so wird $\pm2$ ein biquadratischer Rest von $p$, wenn $g$ von der Form $8n+1$ oder $8n+7$, ein biquadratischer Nichtrest hingegen, wenn $g$ von der Form $8n+3$ oder $8n+5$ ist.

% A544-P0177-D03
Das zweite Criterium hängt zusammen mit der Zerlegung der Zahl $p$ in zwei Quadrate, die bekanntlich auch immer möglich und nur auf Eine Art möglich ist. Setzt man $p=ee+ff$, und nimmt an, dass $ee$ das ungerade, $ff$ das gerade Quadrat bedeutet, so bringt schon die vorausgesetzte Form von $p=8n+1$ mit sich, dass auch $\frac12f$ eine gerade Zahl wird, also $f$ entweder von der Form $8m$ oder von der Form $8m+4$: im ersten Fall nun wird $\pm2$ biquadratischer Rest, im andern biquadratischer Nichtrest von $p$ sein.

% A544-P0177-D04
Wir deuten hier nur die Bemerkung an, wozu die höhere Arithmetik so oft Gelegenheit gibt, dass nicht so wohl die Schönheit und Einfachheit der Theoreme selbst, als die Schwierigkeit ihrer Begründung sie vorzüglich merkwürdig macht. Sobald man einmal veranlasst ist, das Dasein eines Zusammenhanges zwischen dem Verhalten der Zahl $\pm2$ und den beiden angeführten Zerlegungen der Zahl $p$ zu vermuthen, ist es äusserst leicht, diesen Zusammenhang durch Induction wirklich zu entdecken. Allein schon bei dem ersten Criterium ist der Beweis dafür nicht ganz leicht zu führen, viel tiefer versteckt liegt er aber bei dem zweiten, wo er mit anderweitigen subtilen Hülfsuntersuchungen innigst verkettet ist, die ihrerseits wieder zu einer merkwürdigen Erweiterung der Theorie der Kreistheilung führen. Diese wunderbare Verkettung der Wahrheiten ist es vorzüglich. was, wie man schon oft bemerkt hat, der höhern Arithmetik einen so eigenthümlichen Reiz gibt. Diese Begründungen selbst vertragen übrigens natürlich hier keinen Auszug, und müssen in der Abhandlung selbst nachgesehen werden. Allein ein paar andere neue arithmetische Theoreme, welche gleichfalls mit der Begründung des zweiten Criterium innigst verbunden sind, verdienen wohl, ihrer Einfachheit wegen, hier noch besonders herausgehoben zu werden.

% A544-P0177-D05
Wenn $p$ eine Primzahl von der Form $4k+1$ ist, und $=ee+ff$ 
\SourceLeaf{178}%
% A544-P0178-D01
gesetzt wird, so dass $ee$ das ungerade, $ff$ das gerade Quadrat bedeutet; wenn man ferner
\[
1\pdot2\pdot3\mathinner{.\,.\,.\,.\,.\,.}\ k=q
\]
\[
(k+1)(k+2)(k+3)\dotsfive2k=r
\]
setzt, so wird allemal $\pm e$ der kleinste Rest sein, welcher hervorgeht, indem man $\frac r{2q}$ mit $p$ dividirt, und $\pm f$ der kleinste Rest, welchen man aus der Division von $\frac12rr$ mit $p$ erhält (kleinsten Rest immer so verstanden, dass er zwischen den Grenzen $-\frac12p$ und $+\frac12p$ genommen wird). Die Zahl $\frac r{2q}$, welche für $p=5$ den Werth $1$ erhält, kann man für grössere Werthe von $p$ auch in folgende Form setzen
\[
\frac{6\pdot10\pdot14\pdot18\dotsfive(p-3)}{2\pdot3\pdot4\pdot5\dotsfive k}
\]
Es ist sehr merkwürdig, dass so die Zerlegung der Zahl $p$ in zwei Quadrate ganz auf directem Wege erhalten werden kann: aber fast noch merkwürdiger ist ein dabei Statt findender Nebenumstand. Allemal nemlich findet man durch dieses Verfahren die Wurzel des ungeraden Quadrates, $e$, mit positivem Zeichen, wenn $e$, positiv genommen, von der Form $4m+1$ ist, und mit negativem, wenn $e$ positiv genommen von der Form $4m+3$ ist. Hingegen hat für das Zeichen, mit welchem die Wurzel des geraden Quadrats, $f$, aus jener Operation hervorgeht, noch durchaus keine allgemeine Regel aufgefunden werden können, weder a priori, noch auf dem Wege der Induction, und der Verfasser empfiehlt daher, am Schlusse der Abhandlung, diesen Gegenstand den Freunden der höhern Arithmetik zu weiterer Nachforschung, überzeugt, dass mit dem Gelingen derselben sich zugleich eine ergiebige Quelle neuer Erweiterungen dieses schönen Theils der Mathematik eröffnen werde.
\BookRule{.28\linewidth}

\clearpage
\SourceLeaf{179}%
% A544-P0179-D01
\pdfbookmark[0]{Göttingische gelehrte Anzeigen. 1831 April 23.}{notice1831april23}
\begin{center}
\rule{.82\linewidth}{.35pt}\par
Göttingische gelehrte Anzeigen.\quad 1831 April 23.\par
\rule{.82\linewidth}{.35pt}
\end{center}
Eine am 15. April von dem Hofr. \textsc{Gauss} der Königl. Societät überreichte Vorlesung:
\begin{center}\itshape
Theoria residuorum biquadraticorum, commentatio secunda,
\end{center}
ist die Fortsetzung der bereits im sechsten Bande der \emph{Commentationes novae} abgedruckten Abhandlung, wovon auch in unsern Blättern zu seiner Zeit 1825 April 11 eine Anzeige gemacht war. Auch diese Fortsetzung, obgleich mehr als doppelt stärker wie die erste Abhandlung, erschöpft den überaus reichhaltigen Gegenstand noch nicht, und erst einer künftigen dritten Abhandlung wird die Vollendung des Ganzen vorbehalten bleiben.

% A544-P0179-D02
Obgleich die Grundbegriffe dieser Lehren und der Inhalt der ersten Abhandlung als allen, die aus der höhern Arithmetik ein Studium gemacht haben, bekannt vorausgesetzt werden können, wollen wir doch jene zur Bequemlichkeit solcher Freunde dieses Theils der Mathematik, welchen die erste Abhandlung nicht gleich zur Hand ist, hier kurz in Erinnerung bringen. In Beziehung auf eine beliebige ganze Zahl $p$ heisst eine andere $k$ ein biquadratischer Rest, wenn es Zahlen der Form $x^4-k$ gibt, die durch $p$ theilbar sind; im entgegengesetzten Fall heisst sie biquadratischer Nichtrest von $p$. Es ist zureichend, sich hiebei auf den Fall einzuschränken, wo $p$ eine Primzahl der Form $4n+1$, und $k$ durch
\SourceLeaf{180}%
% A544-P0180-D01
dieselbe nicht theilbar ist, da alle andere Fälle entweder für sich klar, oder auf diesen zurückzuführen sind.

% A544-P0180-D02
Für einen solchen \emph{gegebenen} Werth von $p$ zerfallen sämmtliche durch $p$ nicht theilbare Zahlen in \emph{vier} Classen, wovon die eine die biquadratischen Reste, eine zweite solche biquadratische Nichtreste, die quadratische Reste von $p$ sind, enthält, und in die beiden übrigen die biquadratischen Nichtreste, welche zugleich quadratische Nichtreste sind, vertheilt werden. Das Princip dieser Vertheilung besteht darin, dass allemal entweder $k^n-1$, oder $k^n+1$, oder $k^n-f$, oder $k^n+f$ durch $p$ theilbar sein wird, wo $f$ eine ganze Zahl bedeutet, die $ff+1$ durch $p$ theilbar macht. Jeder, dem die elementarische Terminologie bekannt ist, sieht von selbst, wie diese Worterklärungen in dieselbe eingekleidet werden.

% A544-P0180-D03
Die Theorie dieser Classificirung nicht nur für den an der Oberfläche liegenden Fall $k=-1$, sondern auch für die, subtile Hülfsuntersuchungen erfordernden, Fälle $k=\pm2$, findet sich in der ersten Abhandlung ganz vollendet. Im Anfang der gegenwärtigen Abhandlung wird nun zu grösseren Werthen von $k$ fortgeschritten: man braucht aber dabei zunächst nur solche in Betracht zu ziehen, die selbst Primzahlen sind, und der Erfolg zeigt, dass die Resultate am einfachsten ausfallen, wenn man die Werthe positiv oder negativ nimmt, je nachdem sie, absolut betrachtet, von der Form $4m+1$ oder $4m+3$ sind. Die Induction gibt hier sofort mit grosser Leichtigkeit eine reiche Ernte von neuen Lehrsätzen, wovon wir hier nur ein paar anführen. Die Numerirung der Classen mit $1,2,3,4$ wird auf die Fälle bezogen, wo $k^n$ den Zahlen $1,f,-1,-f$ congruent wird; zugleich ist für die Zahl $f$ immer derjenige Werth angenommen, welcher $a+bf$ durch $p$ theilbar macht, wenn $aa+bb$ die Zerlegung von $p$ in ein ungerades und ein gerades Quadrat vorstellt. So findet sich durch die Induction, dass die Zahl $-3$ allemal zu der Classe $1,2,3,4$ gehört, je nachdem $b,a+b,a,a-b$ durch $3$ theilbar ist; dass die Zahl $+5$ der Reihe nach zu jenen Classen gehört, je nachdem $b,a-b,a,a+b$ durch $5$ theilbar ist; dass die Zahl $-7$ in die Classe $1$ fällt, wenn $a$ oder $b$; in die Classe $2$, wenn $a-2b$ oder $a-3b$; in die Classe $3$, wenn $a-b$ oder $a+b$; in die Classe $4$, wenn $a+2b$ oder $a+3b$ durch $7$ theilbar ist. Aehnliche Theoreme ergeben sich in Beziehung auf die Zahlen $-11,+13,+17,-19,-23$ u. s. f. So leicht sich aber alle dergleichen specielle Theoreme durch die Induction entdecken lassen, so schwer scheint
